\begin{lemma}[Projective space bundle formula]
\label{lemma-chow-ring-projective-bundle}
Let $(S, \delta)$ be as in Situation \ref{situation-setup}.
Let $X$ be locally of finite type over $S$.
Let $\mathcal{E}$ be a finite locally free $\mathcal{O}_X$-module
$\mathcal{E}$ of rank $r$. Let $(\pi : P \to X, \mathcal{O}_P(1))$
be the projective bundle associated to $\mathcal{E}$.
The map
$$
\bigoplus\nolimits_{i = 0}^{r - 1}
\CH_{k + i}(X)
\longrightarrow
\CH_{k + r - 1}(P),
$$
$$
(\alpha_0, \ldots, \alpha_{r-1})
\longmapsto
\pi^*\alpha_0 +
c_1(\mathcal{O}_P(1)) \cap \pi^*\alpha_1
+ \ldots +
c_1(\mathcal{O}_P(1))^{r - 1} \cap \pi^*\alpha_{r-1}
$$
is an isomorphism.
\end{lemma}

\begin{proof}
Fix $k \in \mathbf{Z}$. We first show the map is injective.
Suppose that $(\alpha_0, \ldots, \alpha_{r - 1})$ is an element
of the left hand side that maps to zero.
By Lemma \ref{lemma-cap-projective-bundle} we see that
$$
0 = \pi_*(\pi^*\alpha_0 +
c_1(\mathcal{O}_P(1)) \cap \pi^*\alpha_1
+ \ldots +
c_1(\mathcal{O}_P(1))^{r - 1} \cap \pi^*\alpha_{r-1})
= \alpha_{r - 1}
$$
Next, we see that
$$
0 = \pi_*(c_1(\mathcal{O}_P(1)) \cap (\pi^*\alpha_0 +
c_1(\mathcal{O}_P(1)) \cap \pi^*\alpha_1
+ \ldots +
c_1(\mathcal{O}_P(1))^{r - 2} \cap \pi^*\alpha_{r - 2}))
= \alpha_{r - 2}
$$
and so on. Hence the map is injective.

\medskip\noindent
It remains to show the map is surjective.
Let $X_i$, $i \in I$ be the irreducible components of $X$.
Then $P_i = \mathbf{P}(\mathcal{E}|_{X_i})$, $i \in I$
are the irreducible components of $P$. Consider the commutative
diagram
$$
\xymatrix{
\coprod P_i \ar[d]_{\coprod \pi_i} \ar[r]_p & P \ar[d]^\pi \\
\coprod X_i \ar[r]^q & X
}
$$
Observe that $p_*$ is surjective. If $\beta \in \CH_k(\coprod X_i)$
then $\pi^* q_* \beta = p_*(\coprod \pi_i)^* \beta$, see
Lemma \ref{lemma-flat-pullback-proper-pushforward}. Similarly for
capping with $c_1(\mathcal{O}(1))$ by
Lemma \ref{lemma-pushforward-cap-c1}.
Hence, if the map of the lemma is surjective for each
of the morphisms $\pi_i : P_i \to X_i$, then the map is
surjective for $\pi : P \to X$. Hence we may assume $X$ is irreducible.
Thus $\dim_\delta(X) < \infty$ and in particular we may use
induction on $\dim_\delta(X)$.

\medskip\noindent
The result is clear if $\dim_\delta(X) < k$.
Let $\alpha \in \CH_{k + r - 1}(P)$.
For any locally closed subscheme $T \subset X$ denote
$\gamma_T : \bigoplus \CH_{k + i}(T) \to \CH_{k + r - 1}(\pi^{-1}(T))$
the map
$$
\gamma_T(\alpha_0, \ldots, \alpha_{r - 1})
= \pi^*\alpha_0 + \ldots +
c_1(\mathcal{O}_{\pi^{-1}(T)}(1))^{r - 1} \cap \pi^*\alpha_{r - 1}.
$$
Suppose for some nonempty open $U \subset X$ we have
$\alpha|_{\pi^{-1}(U)} = \gamma_U(\alpha_0, \ldots, \alpha_{r - 1})$.
Then we may choose lifts $\alpha'_i \in \CH_{k + i}(X)$ and we
see that $\alpha - \gamma_X(\alpha'_0, \ldots, \alpha'_{r - 1})$
is by Lemma \ref{lemma-restrict-to-open}
rationally equivalent to a $k$-cycle on $P_Y = \mathbf{P}(\mathcal{E}|_Y)$
where $Y = X \setminus U$ as a reduced closed subscheme.
Note that $\dim_\delta(Y) < \dim_\delta(X)$.
By induction the result holds
for $P_Y \to Y$ and hence the result holds for $\alpha$.
Hence we may replace $X$ by any nonempty open of $X$.

\medskip\noindent
In particular we may assume that $\mathcal{E} \cong \mathcal{O}_X^{\oplus r}$.
In this case $\mathbf{P}(\mathcal{E}) = X \times \mathbf{P}^{r - 1}$.
Let us use the stratification
$$
\mathbf{P}^{r - 1} = \mathbf{A}^{r - 1}
\amalg \mathbf{A}^{r - 2}
\amalg \ldots
\amalg \mathbf{A}^0
$$
The closure of each stratum is a $\mathbf{P}^{r - 1 - i}$ which is a
representative of $c_1(\mathcal{O}(1))^i \cap [\mathbf{P}^{r - 1}]$.
Hence $P$ has a similar stratification
$$
P = U^{r - 1} \amalg U^{r - 2} \amalg \ldots \amalg U^0
$$
Let $P^i$ be the closure of $U^i$. Let $\pi^i : P^i \to X$
be the restriction of $\pi$ to $P^i$.
Let $\alpha \in \CH_{k + r - 1}(P)$. By
Lemma \ref{lemma-pullback-affine-fibres-surjective}
we can write $\alpha|_{U^{r - 1}} = \pi^*\alpha_0|_{U^{r - 1}}$
for some $\alpha_0 \in \CH_k(X)$. Hence the difference
$\alpha - \pi^*\alpha_0$ is the image of some
$\alpha' \in \CH_{k + r - 1}(P^{r - 2})$.
By Lemma \ref{lemma-pullback-affine-fibres-surjective}
again we can write
$\alpha'|_{U^{r - 2}} = (\pi^{r - 2})^*\alpha_1|_{U^{r - 2}}$
for some $\alpha_1 \in \CH_{k + 1}(X)$.
By Lemma \ref{lemma-relative-effective-cartier}
we see that the image of $(\pi^{r - 2})^*\alpha_1$
represents $c_1(\mathcal{O}_P(1)) \cap \pi^*\alpha_1$.
We also see that
$\alpha - \pi^*\alpha_0 - c_1(\mathcal{O}_P(1)) \cap \pi^*\alpha_1$
is the image of some $\alpha'' \in \CH_{k + r - 1}(P^{r - 3})$.
And so on.
\end{proof}